% Study sheet: Volume IV, Chapter 6, Policy Training.
% Numbers: \val{...} from training_studysheet_numbers.tex, written by training_studysheet_numbers.py.
\documentclass{studysheet}
\input{training_studysheet_numbers}
\sheetvolume{4}
\sheetchapter{6}{Policy Training}
\sheetsources{Numbers from the chapter's mlsysim scenarios for the cage-door latch on the warehouse mobile manipulator (illustrative); questions marked assumed use new inputs.}

\begin{document}
% ===================================================================== FRONT: LEARN
\frontside
\begin{bigidea}
A learned policy keeps what its training experience rewarded, not which experience that was. Its
training score can rise while a \hl{hidden dependency} on demonstrations, a simulator or a base model
grows. An admitted action crosses the \hl{causal boundary} and no gradient step takes it back, so the
limits of a policy's evidence are written down before it may propose motion.
\end{bigidea}

\begin{leftcol}
\sheetsec{Key ideas}
\begin{ideas}
\idea{Every regime leaves a dependency no score shows.}{Loss and simulated return measure fit to the
  source, so rankings can invert on hardware.}
\idea{Mode averaging is the optimum.}{Squared error picks the mean of two valid swerves at a demonstrated
  state. A multimodal decoder fixes it; more data does not.}
\idea{Drift is a separate failure.}{Errors reach undemonstrated states, which a support monitor flags;
  recovery labels help only where a return is still possible.}
\idea{Slack, not mean latency, decides a decoder.}{Each dependent evaluation adds time. A decoder whose
  slack is below the jitter allowance overruns on one stretch.}
\idea{Randomization covers parameters, not mechanisms.}{Omitted backlash or bus lag stays absent at any
  range; paired probes size each gap.}
\idea{An update is a new checkpoint.}{Old trials describe old weights. Declared ODD = supported cells
  minus omission-log exclusions; other robots add no cell.}
\end{ideas}

\vspace{10pt}
\sheetsec{Five sources of experience}
\begin{sheettable}{@{}l L L@{}}
\thead{Regime} & \thead{Experience} & \thead{Hidden dependency}\\ \midrule
Behavioral cloning & expert demonstrations & coverage: modes and recovery states\\
Offline RL & logged transitions, rewards & support of the logged actions\\
Simulation RL & simulated rollouts & contact model; omitted mechanisms\\
Trial on the machine & real rollouts & calendar, heat, drivetrain wear\\
Pretrain, post-train & other robots' episodes & base checkpoint; forgetting\\
\end{sheettable}
\end{leftcol}%
\begin{rightcol}
\fcard{Decoder latency}%
  {\ProposerInferenceTime = \ObservationEncodingTime + \DecoderEvaluationCount\,\DecoderPassTime \le \CompletionDeadline}%
  {Encode once, then $\DecoderEvaluationCount$ dependent evaluations of $\DecoderPassTime$ each. Synchronous: the deadline is
   the setpoint step, \val{step}\,ms. Asynchronous: the renewal period, \val{renew}\,ms, which must keep a
   \val{jitter}\,ms jitter allowance free. The deployed one-pass ACT decoder takes \val{act}\,ms.}%
  {A fast decoder still needs a lease and the permission check.}{}
\fcard{Trial campaign on the machine}%
  {\TrialCampaignWallTime = \frac{\TrialCount\,(\TrialRolloutTime + \TrialResetTime)}{\ParallelRigCount},
   \qquad \text{wear} = \TrialCount\,\TrialLoadReversalCount}%
  {$\TrialCount$ trials of rollout $\TrialRolloutTime$ and reset $\TrialResetTime$ over $\ParallelRigCount$ rigs;
   $\TrialLoadReversalCount$ drivetrain load reversals per trial. Rigs divide the calendar, never the wear.}%
  {Latch: \val{trials} trials, \val{rollout}\,s each, \val{revper} reversals; \val{daysmanual} days on one rig with
   \val{resetmanual}\,s manual resets.}{}
\fcard{Impact peak}%
  {\UnbrakedImpactForce = \SpeedAtDetection\sqrt{\ContactStiffness\,\Mass}}%
  {A mass $\Mass$ coasting at $\SpeedAtDetection$ into contact stiffness $\ContactStiffness$. A soft simulated spring
   understates the real plate's peak, so an unbraked plunge scores well in training.}%
  {Latch: $\Mass$ = \val{mass}\,kg, plate \val{kplate}\,N/m, penalty spring \val{ksim}\,N/m, limit \val{flim}\,N.
   The plate alone is the rigid-arm bound; arm compliance lowers it, so a push test must measure the arm.}%
  {\begin{tikzpicture}[sheet]
     \pgfmathsetmacro\sc{0.85}
     \draw[nbox] (70,22) rectangle ({70+\val{bsim}*\sc},32);
     \draw[hotbox] (70,6) rectangle ({70+\val{brigid}*\sc},16);
     \node[lab, anchor=east] at (66,27) {simulated spring};
     \node[lab, anchor=east] at (66,11) {real plate};
     \node[note, anchor=west] at ({73+\val{bsim}*\sc},27) {\val{fsim}\,N};
     \node[note, anchor=west] at ({73+\val{brigid}*\sc},11) {\val{frigid}\,N};
     \draw[dim, dashed] ({70+\val{blim}*\sc},0) -- ({70+\val{blim}*\sc},40);
     \draw[white, dashed, line width=0.6pt] ({70+\val{blim}*\sc},6) -- ({70+\val{blim}*\sc},16);
     \node[small, anchor=south] at ({70+\val{blim}*\sc},40) {latch limit \val{flim}\,N};
     \node[note, anchor=north west] at (70,-2) {peak force at \val{vfast}\,m/s, rigid arm};
   \end{tikzpicture}}
\end{rightcol}

% ===================================================================== BACK: PRACTICE
\backside
\begin{answers}
\ans{Finish}{\val{turnlat}\,ms $>$ \val{step}\,ms: still not synchronous. Slack \val{renew} $-$ \val{turnlat} =
  \res{\val{turnslack}\,ms}, above the \val{jitter}\,ms allowance; a stretch ends at \val{turnend}\,ms, inside the
  renewal: asynchronous under a lease, if action quality holds (decoder card).}
\ans{1}{\val{trials} $\times$ (\val{rollout} + \val{qreset})\,s $\div$ \val{rigs} = \val{qsec}\,s $\approx$ \res{\val{qdays}\,days}.
  Wear is still \res{\val{reversals} reversals}: rigs divide the calendar, not the drivetrain's load (campaign card).}
\ans{2}{Plate: \val{qv} $\times$ \val{qroot} = \res{\val{qfrigid}\,N}, \val{qover}\,N over the limit. Simulator: \val{qv} $\times$
  \val{qrootsim} = \val{qfsim}\,N, looks safe. The plate counts until a push test measures the arm's stiffness (impact card).}
\columnbreak
\ans{3}{Column: mode averaging at a demonstrated observation with two valid modes; fix with a multimodal
  decoder. Sagged door: drift into undemonstrated states, which a support monitor flags; fix with recovery
  labels (key ideas on modes and drift).}
\ans{4}{\res{\val{supported} of \val{cells} cells}. Other machines' episodes join no latch cell; plausible moves at the displaced plate are borrowed competence, not evidence (key idea on checkpoints).}
\end{answers}

\begin{leftcol}
\sheetsec{Worked example}
\begin{practicebox}\boxtitle{Can a diffusion head replace the one-pass decoder?}
Multimodal demonstrations tempt a swap of ACT (\val{act}\,ms) for a \val{evals}-step diffusion head:
encoding \val{enc}\,ms, total \val{ddim}\,ms. Setpoint step \val{step}\,ms, renewal \val{renew}\,ms, jitter
allowance \val{jitter}\,ms.\par\vspace{5pt}
{\centering
\begin{tikzpicture}[sheet]
  \draw[nbox] (40,30) rectangle (40+\val{pact},40);
  \draw[accbox] (40,12) rectangle (40+\val{penc},22);
  \draw[hotbox] (40+\val{penc},12) rectangle (40+\val{pddim},22);
  \draw[accbox, dashed] (40+\val{pddim},12) rectangle (40+\val{pjit},22);
  \node[lab, anchor=east] at (36,35) {ACT};
  \node[lab, anchor=east] at (36,17) {diffusion};
  \node[small] at ({40+\val{penc}/2},6) {encode};
  \draw[dim, dashed] (40+\val{pstep},0) -- (40+\val{pstep},48);
  \draw[white, dashed, line width=0.6pt] (40+\val{pstep},12) -- (40+\val{pstep},22);
  \draw[dim, dashed] (40+\val{prenew},0) -- (40+\val{prenew},48);
  \node[small, anchor=south] at (40+\val{pstep},48) {setpoint step};
  \node[small, anchor=south] at (40+\val{prenew},48) {renewal};
  \node[note, anchor=north] at (40+\val{prenew}/2,-3) {to scale; dashed box = jitter stretch};
\end{tikzpicture}\par}\vspace{4pt}
\begin{steps}
\item Per evaluation: (\val{ddim} $-$ \val{enc})\,ms $\div$ \val{evals} = \res{\val{pereval}\,ms}.
\item Synchronous: \val{ddim}\,ms $>$ \val{step}\,ms, so it misses every setpoint.
\item Renewal slack: \val{renew} $-$ \val{ddim} = \res{\val{slackddim}\,ms} (ACT: \val{slackact}\,ms).
\item Stretched: \val{ddim} + \val{jitter} $-$ \val{renew} = \res{\val{overrun}\,ms} past renewal; ACT ends at \val{actend}\,ms, inside.
\end{steps}
\takeaway{Asynchronous only, and one stretch overruns: the old prefix plays on until its lease expires, then the fallback acts.}
\end{practicebox}

\vspace{10pt}
\sheetsec{Your turn: finish it}
\begin{practicebox}\boxtitle{Same head, distilled to \val{turnevals} evaluations}
\assumed{} step distillation cuts the head to \val{turnevals} evaluations; other costs are unchanged.\par\vspace{4pt}
\textbf{Given:} $\ProposerInferenceTime$ = \val{enc} + \val{turnevals} $\times$ \val{pereval} = \val{turnlat}\,ms.\par
\textbf{You:} does it fit the setpoint step? What renewal slack remains, and does a jitter stretch still overrun?\par\vspace{30pt}
\end{practicebox}
\end{leftcol}%
\begin{rightcol}
\sheetsec{Common mistakes}
\mistake{``Our regression policy has the lowest validation loss, so it won't steer into the column.''}%
  {The mean minimizes squared error on held-out data from the same two modes. Swerves at plus and minus
   yaw rate average to zero yaw: the best score, aimed at the column.}{}
\mistake{``Widen the randomization ranges and the simulator's gaps close.''}%
  {Randomization varies masses and friction already in the equations. Backlash and bus lag stay absent at
   any range; they go in the omission log until paired probes bound them.}{}
\mistake{``Fine-tuning fixed the fixture error, so the new weights are ready.''}%
  {Earlier trials describe the old weights only; free-space and stopping behavior may regress. Limit the
   correction to contact, then retest contact, free motion and stopping.}{}

\vspace{2pt}
\sheetsec{Check yourself}
\begin{qs}
\item \assumed{} a faster fixture resets in \val{qreset}\,s, and the latch campaign (\val{trials} trials of
  \val{rollout}\,s, \val{revper} reversals each) runs on \val{rigs} rigs. Calendar days? Total load reversals?
\item \assumed{} the policy coasts into the plate at \val{qv}\,m/s ($\Mass$ = \val{mass}\,kg). Peak force against
  the real plate (rigid arm) and against the penalty spring? Which do you trust?
\item One cloned policy steers at a column it was taught to pass; another binds on a sagged door. Name
  each failure, the remedy, and the observation that tells them apart.
\item The latch demonstrations support \val{supported} of \val{cells} scenario cells. A team post-trains the
  \val{basename} base on them, and it moves plausibly at the displaced plate. How many cells does its declared ODD hold?
\end{qs}
\end{rightcol}
\end{document}
